\documentclass[../../main.tex]{subfiles}

\begin{document}

\section{Smirnov 度量化定理}\label{section:ch06-smirnov-metrization}

Nagata-Smirnov 度量化定理给出了空间可度量化的一个充分必要条件. 在本节中, 我们给出另一个度量化定理, 它是 Nagata-Smirnov 定理的一个推论, 是由 Smirnov 首先证明的.

\begin{definition}\label{definition:ch06-locally-metrizable}
称空间 $X$ 是局部可度量化的(locally metrizable), 如果 $X$ 的每一点 $x$ 有一个邻域 $U$ 在子空间拓扑之下是可度量化的.
\end{definition}

\begin{theorem}[Smirnov 度量化定理]\label{theorem:ch06-smirnov-metrization}
空间 $X$ 可度量化当且仅当 $X$ 是仿紧致的 Hausdorff 空间, 并且是局部可度量化的.
\end{theorem}

\begin{proof}
假定 $X$ 可度量化. 根据\cref{theorem:ch06-metrizable-paracompact}可见, $X$ 是局部可度量化的, 也是仿紧致的.

反之, 假设 $X$ 是仿紧致的 Hausdorff 空间, 并且是局部可度量化的. 我们只需证明 $X$ 有一个可数局部有限基. 因为 $X$ 是正则的, 根据\cref{theorem:ch06-nagata-smirnov-metrization}可见 $X$ 是可度量化的.

将\cref{theorem:ch06-nagata-smirnov-metrization}证明的后一部分调整一下就是这里的证明. 先用一些可度量化的开集来覆盖 $X$, 然后选取这个覆盖的一个局部有限开加细 $\mathcal{C}$ 覆盖 $X$. $\mathcal{C}$ 的每一个成员 $C$ 都是可度量化的. 设函数 $d_C\colon C\times C\to\mathbb{R}$ 是给出 $C$ 的拓扑的一个度量. 给定 $x\in X$, 用 $B_C(x,\varepsilon)$ 表示 $C$ 中所有满足 $d_C(x,y)<\varepsilon$ 的点 $y$ 构成的集合, 因为 $B_C(x,\varepsilon)$ 是 $C$ 的一个开集, 所以也是 $X$ 的一个开集.

给定 $m\in\mathbb{Z}_+$, 令 $\mathcal{A}_m$ 为半径是 $1/m$ 的开球构成的 $X$ 的覆盖, 即
\begin{align*}
\mathcal{A}_m=\{B_C(x,1/m)\mid x\in C\text{ 且 }C\in\mathcal{C}\}.
\end{align*}
设 $\mathcal{D}_m$ 为 $\mathcal{A}_m$ 的覆盖 $X$ 的一个局部有限开加细. (此处用到了仿紧致性.) 设 $\mathcal{D}$ 是这些集族 $\mathcal{D}_m$ 的并. 则 $\mathcal{D}$ 是可数局部有限的. 我们断言 $\mathcal{D}$ 是 $X$ 的一个基, 从而定理得证.

设 $x$ 是 $X$ 的一个点, $U$ 是 $x$ 的一个邻域. 我们设法找到 $\mathcal{D}$ 的一个成员 $D$, 使得 $x\in D\subset U$. 现在, $x$ 仅属于 $\mathcal{C}$ 中有限多个成员, 比如说属于 $C_1,\cdots,C_k$. 于是在集合 $C_i$ 中, $U\cap C_i$ 是 $x$ 的一个邻域, 所以存在一个 $\varepsilon_i>0$ 使得
\begin{align*}
B_{C_i}(x,\varepsilon_i)\subset(U\cap C_i).
\end{align*}
选取 $m$, 使得 $2/m<\min\{\varepsilon_1,\cdots,\varepsilon_k\}$. 因为族 $\mathcal{D}_m$ 覆盖 $X$, 因此 $\mathcal{D}_m$ 中必存在一个成员 $D$ 包含 $x$. 因为 $\mathcal{D}_m$ 加细 $\mathcal{A}_m$, 所以对某一个 $C\in\mathcal{C}$ 及某一个 $y\in C$, 必存在 $\mathcal{A}_m$ 的一个成员 $B_C(y,1/m)$ 包含 $D$. 因为
\begin{align*}
x\in D\subset B_C(y,1/m),
\end{align*}
所以 $x$ 属于 $C$, 从而 $C$ 必然是集合 $C_1,\cdots,C_k$ 中的一个, 比如说 $C=C_i$. 又因 $B_C(y,1/m)$ 的直径最多为 $2/m<\varepsilon_i$, 于是, 我们有
\begin{align*}
x\in D\subset B_{C_i}(y,1/m)\subset B_{C_i}(x,\varepsilon_i)\subset U.
\end{align*}
\end{proof}

\subsection*{习题}

\begin{enumerate}
\item 比较\cref{theorem:ch06-smirnov-metrization}和第 34 节的习题 7 和习题 8.
\item (a) 证明: 对于每个 $x\in S_\Omega$, 由 $x$ 确定的 $S_\Omega$ 的截有可数基, 因此是可度量化的.
\begin{enumerate}[(b)]
\item 证明 $S_\Omega$ 不是仿紧致的.
\end{enumerate}
\end{enumerate}

\end{document}